Large language models can be manipulated by adversarial prompts that appear as gibberish to humans, yet no prior work has systematically tested whether models can \emph{explain} the nonsense prompts they successfully follow.
We present the first study of LLM self-explanation ability across a controlled perplexity spectrum, evaluating \gptfour on 40~prompts spanning five categories: natural language, obfuscated text, human-crafted jailbreaks, \gcg adversarial suffixes, and random token strings.
We find a striking U-shaped pattern in explanation quality: models explain natural language well (4.5/5) and \emph{perfectly} identify pure random tokens as meaningless (5.0/5), but struggle most with partially structured nonsense---\gcg suffixes and obfuscated text score only 3.6--3.75/5.
This \uncannyvalley of prompt comprehension arises because models possess two distinct explanation mechanisms: genuine \emph{understanding} for natural language and \metarecog for pure noise, with a gap between them that adversarial attacks exploit.
We further show that 5 of 8 \gcg suffixes trigger compliant responses with confabulated justifications, and that random search for directed high-perplexity prompts succeeds in only 5\% of trials.
These findings reveal that the hardest prompts for LLMs to interpret are not the most nonsensical, but those with just enough structure to trigger confabulated interpretations---precisely the regime where safety training fails (Kruskal-Wallis $H = 14.96$, $p = 0.005$).
